\usection{Antecedentes de la Investigación}

\usubsection{El desarrollo de Scratch-Jr: el aprendizaje de programación en primera infancia como nueva alfabetización}

En este trabajo \cite{bers2023} se centra en el desarrollo de ScratchJr junto a Mitchel Resnick, examina cómo el aprendizaje de la programación visual y por bloques puede funcionar como un “nueva alfabetización” para niños de 5 a 7 años. En el artículo describe el enfoque pedagógico de “código como otro lenguaje”, que concibe la enseñanza de la programación no solo como una herramienta para resolver problemas, sino como una forma de expresión personal que permite que los niños creen, compartan y desarrollen ideas para aprender socializando.

\usubsection{Entorno de robótica educativa multiagente orientado a favorecer el desarrollo del pensamiento computacional en jóvenes cursantes de educación media}

El trabajo de \cite{espejo2022} se enfoca en el desarrollo de un entorno de robótica multiagente diseñado para fomentar el pensamiento computacional en jóvenes de educación media. Este proyecto es particularmente relevante para la presente investigación porque demuestra la efectividad de los entornos tangibles y visuales para enseñar conceptos complejos de concurrencia y colaboración.

En específico, Espejo implementó un sistema que permite a los estudiantes programar agentes robóticos para que interactúen en un espacio de realidad aumentada, combinando la programación en bloques de MakeCode con una arquitectura descentralizada para la inteligencia de los agentes. Este enfoque valida la idea de que la colaboración entre múltiples entidades puede ser una poderosa herramienta pedagógica. A diferencia de un sistema de un solo agente, el entorno multiagente de Espejo obliga a los estudiantes a considerar cómo los diferentes componentes deben operar simultáneamente para alcanzar una meta común, lo cual es fundamental para la comprensión de la concurrencia.

\usubsection{Extensiones de programación de aplicaciones distribuidas para fomentar pensamiento computacional en niños de educación básica}

Rodríguez (2024) en su trabajo propone una solución para que los niños de educación básica se relacionen con conceptos más avanzados de la computación, como lo son la concurrencia y los sistemas distribuidos. Esto lo logró mediante la creación de extensiones en el sistema de MakeCode, demostrando que es posible hacer accesible estos conceptos a un público infantil, despertando el interés y validando que los niños puedan comprender y aplicar nociones de sistemas distribuidos. Su modelo de solución propuesta se basó en “cuadernos” que actúan como una memoria compartida. Estos cuadernos permiten a los estudiantes almacenar y acceder a datos de manera remota, facilitando la comprensión de cómo los diferentes dispositivos en una red pueden interactuar y compartir información. Este trabajo aporta elementos teóricos y prácticos sobre la relación entre la programación distribuida y el pensamiento computacional, así como sobre el diseño e implementación de un entorno de programación visual para fomentar conceptos avanzados como la concurrencia en niños de educación básica.

\usubsection{Scratch: programming for all [Scratch: programación para todos]}

En este trabajo, \cite{resnicketal2009} describen las bases y principios tras el diseño de Scratch, un lenguaje de programación visual basado en bloques y flujos de instrucciones que permite a niños y jóvenes aprender conceptos matemáticos y computacionales, además de ejercitar el pensamiento creativo, racional y sistemático, y el trabajo colaborativo. Se apoya en el pensamiento de Papert para diseñar un lenguaje con un piso bajo que representa la facilidad para aprenderlo, un techo alto que representa el potencial creciente de los problemas que permite resolver, y muros amplios que representan la variedad de proyectos que pueden hacerse con Scratch. El diseño se basa en que el lenguaje sea \emph{tinkerable} (modificable, fácil de “jugar con él”), significativo y social.

El trabajo del equipo detrás de Scratch destaca a través de experiencias prácticas cómo un lenguaje de programación visual ha ayudado a niños y jóvenes a explotar su curiosidad y creatividad, formando grupos de trabajo en los que se desarrolla el pensamiento computacional a través de la creación de juegos, animaciones y otras creaciones interactivas y colaborativas dentro del entorno de Scratch. Estos principios de diseño y ejemplos prácticos fortalecen las bases a utilizar en el desarrollo del presente proyecto.